% GENERATED FILE. Edit canonical lessons, then run npm run generate:books.
% canonical-source-hash: fnv1a64:200952405bcf0f1b
% canonical-lessons: SA-C483-gutika, SA-C483-curnam, SA-C483-lepah, SA-C483-kasayah, SA-C483-limpati

\chapter{Medicines}
\label{ch:sa-medicines}
\hlchaptermodality{483}

\begin{chapteropening}
\textbf{By the end of this chapter:} \emph{I can say how a medicine is taken or applied.}

Complete the last lesson of chapter 483: I can say how a medicine is taken or applied.
\end{chapteropening}

\section[\texorpdfstring{\sk{गुटिका} (guṭikā)}{गुटिका (guṭikā)}]{\sk{गुटिका} (guṭikā) — a pill}

\label{lesson:SA-C483-gutika}



\begin{quote}
Before the new one: say the Sanskrit for medical treatment, then the Sanskrit for a splinter.
\end{quote}

\subsection*{You'll want to know: \sk{गुटिका}}
\textbf{\sk{गुटिका}} — \emph{guṭikā} — \textquotedblleft{}a pill\textquotedblright{}.

\begin{quote}
\sk{प्रतिदिनं} \sk{द्वे} \sk{गुटिके} \sk{खादतु।}
\end{quote}

\begin{quote}
\emph{pratidinaṃ dve guṭike khādatu.}
\end{quote}

\begin{quote}
Take two pills every day.
\end{quote}

A feminine noun.

\begin{cousinweb}
\textbf{\sk{गुटिका}} (\emph{guṭikā}) MW's first sense is \textquotedblleft{}a small globe or ball\textquotedblright{} and then \textquotedblleft{}a pill\textquotedblright{}; MW gives no derivation, so its origin is uncertain (MW).
\end{cousinweb}

\begin{grammarlens}[title={What you've built}]
One more word for describing what you lived and felt.
\end{grammarlens}

\subsection*{Your turn}
\begin{itemize}
\raggedright
  \item \emph{Say it:} \emph{guṭikā}
  \item \emph{Say it:} \emph{guṭikā}, once more
  \item \emph{Say it:} \emph{pratidinaṃ dve guṭike khādatu.} — Take two pills every day.
  \item \emph{Recall:} say the Sanskrit for medical treatment, then the Sanskrit for a splinter, then say \emph{guṭikā} again
\end{itemize}

\begin{tcolorbox}[breakable,colback=teal!4,colframe=teal!35!black,title={Before you move on}]
What does \sk{गुटिका} mean? (\textquotedblleft{}A pill\textquotedblright{}.) Say it once more.
\end{tcolorbox}
\section[\texorpdfstring{\sk{चूर्णम्} (cūrṇam)}{चूर्णम् (cūrṇam)}]{\sk{चूर्णम्} (cūrṇam) — a powder}

\label{lesson:SA-C483-curnam}



\begin{quote}
Before the new one: say the Sanskrit for a splinter, then the Sanskrit for a pill.
\end{quote}

\subsection*{You'll want to know: \sk{चूर्णम्}}
\textbf{\sk{चूर्णम्}} — \emph{cūrṇam} — \textquotedblleft{}a powder\textquotedblright{}.

\begin{quote}
\sk{वैद्यः} \sk{उदरपीडायै} \sk{चूर्णं} \sk{ददाति।}
\end{quote}

\begin{quote}
\emph{vaidyaḥ udarapīḍāyai cūrṇaṃ dadāti.}
\end{quote}

\begin{quote}
The physician gives a powder for the stomach ache.
\end{quote}

A neuter noun.

\begin{cousinweb}
\textbf{\sk{चूर्णम्}} (\emph{cūrṇam}) is traced to the root √carv, \textquotedblleft{}to grind with the teeth\textquotedblright{}, by MW, which gives \textquotedblleft{}powder, flour\textquotedblright{} (MW).
\end{cousinweb}

\begin{grammarlens}[title={What you've built}]
One more word for describing what you lived and felt.
\end{grammarlens}

\subsection*{Your turn}
\begin{itemize}
\raggedright
  \item \emph{Say it:} \emph{cūrṇam}
  \item \emph{Say it:} \emph{cūrṇam}, once more
  \item \emph{Say it:} \emph{vaidyaḥ udarapīḍāyai cūrṇaṃ dadāti.} — The physician gives a powder for the stomach ache.
  \item \emph{Recall:} say the Sanskrit for a splinter, then the Sanskrit for a pill, then say \emph{cūrṇam} again
\end{itemize}

\begin{tcolorbox}[breakable,colback=teal!4,colframe=teal!35!black,title={Before you move on}]
What does \sk{चूर्णम्} mean? (\textquotedblleft{}A powder\textquotedblright{}.) Say it once more.
\end{tcolorbox}
\section[\texorpdfstring{\sk{लेपः} (lepaḥ)}{लेपः (lepaḥ)}]{\sk{लेपः} (lepaḥ) — an ointment}

\label{lesson:SA-C483-lepah}



\begin{quote}
Before the new one: say the Sanskrit for a pill, then the Sanskrit for a powder.
\end{quote}

\subsection*{You'll want to know: \sk{लेपः}}
\textbf{\sk{लेपः}} — \emph{lepaḥ} — \textquotedblleft{}an ointment\textquotedblright{}.

\begin{quote}
\sk{माता} \sk{व्रणे} \sk{लेपं} \sk{स्थापयति।}
\end{quote}

\begin{quote}
\emph{mātā vraṇe lepaṃ sthāpayati.}
\end{quote}

\begin{quote}
Mother puts ointment on the wound.
\end{quote}

A masculine noun.

\begin{cousinweb}
\textbf{\sk{लेपः}} (\emph{lepaḥ}) comes from the root √lip, \textquotedblleft{}to smear\textquotedblright{}; MW gives \textquotedblleft{}smearing\textquotedblright{} and \textquotedblleft{}ointment, unguent, plaster\textquotedblright{} (MW).
\end{cousinweb}

\begin{grammarlens}[title={What you've built}]
One more word for describing what you lived and felt.
\end{grammarlens}

\subsection*{Your turn}
\begin{itemize}
\raggedright
  \item \emph{Say it:} \emph{lepaḥ}
  \item \emph{Say it:} \emph{lepaḥ}, once more
  \item \emph{Say it:} \emph{mātā vraṇe lepaṃ sthāpayati.} — Mother puts ointment on the wound.
  \item \emph{Recall:} say the Sanskrit for a pill, then the Sanskrit for a powder, then say \emph{lepaḥ} again
\end{itemize}

\begin{tcolorbox}[breakable,colback=teal!4,colframe=teal!35!black,title={Before you move on}]
What does \sk{लेपः} mean? (\textquotedblleft{}An ointment\textquotedblright{}.) Say it once more.
\end{tcolorbox}
\section[\texorpdfstring{\sk{कषायः} (kaṣāyaḥ)}{कषायः (kaṣāyaḥ)}]{\sk{कषायः} (kaṣāyaḥ) — a decoction}

\label{lesson:SA-C483-kasayah}



\begin{quote}
Before the new one: say the Sanskrit for a powder, then the Sanskrit for an ointment.
\end{quote}

\subsection*{You'll want to know: \sk{कषायः}}
\textbf{\sk{कषायः}} — \emph{kaṣāyaḥ} — \textquotedblleft{}a decoction\textquotedblright{}.

\begin{quote}
\sk{पितामही} \sk{आर्द्रकस्य} \sk{कषायं} \sk{पचति।}
\end{quote}

\begin{quote}
\emph{pitāmahī ārdrakasya kaṣāyaṃ pacati.}
\end{quote}

\begin{quote}
Grandmother boils a ginger decoction.
\end{quote}

A masculine noun.

\begin{cousinweb}
\textbf{\sk{कषायः}} (\emph{kaṣāyaḥ}) MW gives \textquotedblleft{}astringent\textquotedblright{} and, in medicine, \textquotedblleft{}a decoction\textquotedblright{} boiled down to a quarter; MW gives no derivation, so its origin is uncertain (MW).
\end{cousinweb}

\begin{grammarlens}[title={What you've built}]
One more word for describing what you lived and felt.
\end{grammarlens}

\subsection*{Your turn}
\begin{itemize}
\raggedright
  \item \emph{Say it:} \emph{kaṣāyaḥ}
  \item \emph{Say it:} \emph{kaṣāyaḥ}, once more
  \item \emph{Say it:} \emph{pitāmahī ārdrakasya kaṣāyaṃ pacati.} — Grandmother boils a ginger decoction.
  \item \emph{Recall:} say the Sanskrit for a powder, then the Sanskrit for an ointment, then say \emph{kaṣāyaḥ} again
\end{itemize}

\begin{tcolorbox}[breakable,colback=teal!4,colframe=teal!35!black,title={Before you move on}]
What does \sk{कषायः} mean? (\textquotedblleft{}A decoction\textquotedblright{}.) Say it once more.
\end{tcolorbox}
\section[\texorpdfstring{\sk{लिम्पति} (limpati)}{लिम्पति (limpati)}]{\sk{लिम्पति} (limpati) — smears}

\label{lesson:SA-C483-limpati}



\begin{quote}
Before the new one: say the Sanskrit for an ointment, then the Sanskrit for a decoction.
\end{quote}

\subsection*{You'll want to know: \sk{लिम्पति}}
\textbf{\sk{लिम्पति}} — \emph{limpati} — \textquotedblleft{}smears\textquotedblright{}.

\begin{quote}
\sk{वैद्यः} \sk{व्रणे} \sk{लेपं} \sk{लिम्पति।}
\end{quote}

\begin{quote}
\emph{vaidyaḥ vraṇe lepaṃ limpati.}
\end{quote}

\begin{quote}
The physician smears ointment on the wound.
\end{quote}

Past (he …): \sk{सः} \sk{व्रणे} \sk{लेपं} \sk{लिप्तवान्}, \emph{saḥ vraṇe lepaṃ liptavān}.

\begin{cousinweb}
\textbf{\sk{लिम्पति}} (\emph{limpati}) comes from the root √lip, \textquotedblleft{}to smear, to anoint\textquotedblright{}, whose present stem takes a nasal (\emph{li-m-pa}) (MW).
\end{cousinweb}

\begin{grammarlens}[title={What you've built}]
One more word for describing what you lived and felt.
\end{grammarlens}

\subsection*{Your turn}
\begin{itemize}
\raggedright
  \item \emph{Say it:} \emph{limpati}
  \item \emph{Say it:} \emph{limpati}, once more
  \item \emph{Say it:} \emph{vaidyaḥ vraṇe lepaṃ limpati.} — The physician smears ointment on the wound.
  \item \emph{Recall:} say the Sanskrit for an ointment, then the Sanskrit for a decoction, then say \emph{limpati} again
\end{itemize}

\begin{tcolorbox}[breakable,colback=teal!4,colframe=teal!35!black,title={Before you move on}]
What does \sk{लिम्पति} mean? (\textquotedblleft{}Smears\textquotedblright{}.) Say it once more.
\end{tcolorbox}
